\section*{Practical Appendices}
\addcontentsline{toc}{section}{Practical Appendices}

\subsection{Appendix A: Founder Corporate Facts Audit}

\begin{itemize}[leftmargin=*,itemsep=1.5pt]
\item[{[ ]}] Legal company name, CIN, registered office address, and statutory business objects.
\item[{[ ]}] Directors, designated signatories, and DIN verification.
\item[{[ ]}] Current issued, paid-up, and fully diluted ownership breakdown.
\item[{[ ]}] Shareholder agreements, articles of association, convertible rights, and ESOPs.
\item[{[ ]}] Current tax (PAN, TAN, GSTIN) and payroll registrations (PF, ESIC, PT).
\item[{[ ]}] Statutory filing and annual audit status with the Registrar of Companies (ROC).
\item[{[ ]}] Bank signing authority, mandate limits, and internet banking dual-authorization rules.
\item[{[ ]}] Brand, domain, trademark, curriculum, and software repository ownership.
\item[{[ ]}] Employment agreements and prior-employer outside-business/invention conflict reviews.
\item[{[ ]}] Material customer, vendor, employment, contractor, and IP assignment contracts.
\item[{[ ]}] Privacy and data-processing inventory under Indian DPDP Act and 2025 Rules.
\item[{[ ]}] Current or contemplated debt instruments, director personal guarantees, and asset charges.
\item[{[ ]}] Existing litigation, tax demands, show-cause notices, related parties, and formal disputes.
\item[{[ ]}] Actual legal status of the Bayes brand versus the registered corporate operating entity.
\item[{[ ]}] Independent findings from company secretary, legal counsel, and statutory auditor.
\end{itemize}

\begin{table*}[t]
\centering
\small
\caption{Investor Term Sheet Negotiation Grid: Balancing Investor Protection and Founder Agency}
\label{tab:ch35_term_sheet_grid}
\begin{tabularx}{\textwidth}{p{3.8cm} p{3.2cm} p{3.5cm} X}
\toprule
\textbf{Term} & \textbf{Investor Purpose} & \textbf{Founder Concern} & \textbf{Negotiation Question} \\
\midrule
Valuation and option pool & Economics and hiring & Effective dilution & Before or after financing? \\
Liquidation preference & Downside protection & Low-exit waterfall & Participating? Senior? Cap? \\
Board seats/quorum & Oversight & Decision blockage & Alternate quorum process? \\
Reserved matters & Protect principal & Operational veto & Thresholds, sunset, response time? \\
Anti-dilution & Down-round protection & Transfer of equity & Broad/narrow basis; carve-outs? \\
Drag-along & Enable exit & Forced sale & Trigger and price protection? \\
Founder vesting/lock-in & Continuity & Involuntary forfeiture & Good leaver, disability, cure? \\
Information rights & Transparency & Customer/data confidentiality & Aggregated reports and permitted access? \\
Redemption/debt covenants & Liquidity & Forced cash obligation & Can company survive adverse case? \\
\bottomrule
\end{tabularx}
\end{table*}

\subsection{Appendix B: Investor Term Sheet Negotiation Grid}

Table~\ref{tab:ch35_term_sheet_grid} outlines the core negotiation vectors encountered by founders when raising institutional venture capital or growth equity. Each clause must be evaluated not in isolation, but as an interdependent package of economic returns, negative control, and operational agency.

\subsection{Appendix C: Fully Diluted Cap Table Worksheet}

Track each class and instrument: issued common, preferred, options granted and available pool, warrants, SAFEs/notes, conversion discounts/caps, anti-dilution adjustment and accrued rights. Show pre-money, new money, post-money, owner percentages, proceeds and control rights. Run dilution and exit waterfalls at several valuations including disappointing exits.

\subsection{Appendix D: Reserved Matters Authority Schedule}

Rows: ordinary contracts, discount/refund, hiring, compensation, borrowing, new equity, capex, related-party dealings, material IP licenses, acquisitions, land, research grants, public educational claims and changes to academic policy. Columns: manager authority, CEO authority, board resolution, investor consent, statutory/member approval, evidence and response SLA. Counsel must validate the final schedule.

\subsection{Appendix E: Corporate Compliance Register}

Entity; obligation and legal basis; regulator; applicability threshold; statutory due date; operational due date; responsible internal owner; external adviser; filing evidence; penalty/impact; outstanding issue. Separate company, tax, labor, data, foreign investment, premises and academic requirements.

\subsection{Appendix F: Related-Party Transaction Memo}

Parties and ownership links; commercial rationale; alternatives; independent market quotes; valuation; service specification; conflicted decision makers; disclosures and recusals; approvals under laws/articles; accounting; contract term; termination; donor/academic implications; future review date.

\subsection{Appendix G: Academic Quality and Integrity Charter}

Publish competency standards and assessment rules; faculty academic discretion bounded by policy; 2nd-reader appeal and integrity procedures; truthful outcomes; research ethics and authorship; learner and staff grievance rights; sponsor publication rules; independent review for conflicts and serious misconduct.

\subsection{Appendix H: Philanthropy and Endowment Separation}

Record donor, legal entity, agreement, date/currency, purpose restriction, release conditions, bank location, spend authority, investment policy, reporting, amount collected versus pledged, funds used, balance and contingencies. Never treat restricted gifts as general operating cash without legal basis.

\subsection{Appendix I: Governance Risk Map}

Rank scenarios by legal severity, cash impact, likelihood range, detection delay, control effectiveness and owner: investor veto deadlock; statutory filing missed; delayed college collections; misleading credential claims; data breach; related-party lease; research misconduct; unauthorized foreign grant; founder incapacitation; campus cost overrun. Require quantified triggers and a response plan.

\subsection{Appendix J: A 90-Day Founder Governance Upgrade}

\textbf{Days 1-15:} reconcile entities, shares, banks and registrations. \textbf{16-30:} contract/IP and compliance inventory. \textbf{31-45:} cash controls and delegation thresholds. \textbf{46-60:} privacy, refunds, grievances and employment agreements. \textbf{61-75:} investor/data room readiness and related-party register. \textbf{76-90:} external legal/accounting review, risk register and board-style decision cadence. Complete the critical legal items before pursuing large deals.

\subsection{Appendix K: Shareholder and Academic Control Separation}

For each domain, name the authority and protection: shareholder economic rights; directors' fiduciary duties; company operating budgets; faculty appointment process; academic standards; grant spending; research publications; student admissions and appeals; university sponsor functions; philanthropy; property rights; and IP licensing. Mark any potential conflicts requiring outside review.

\subsection{Appendix L: Founder Incapacity and Succession File}

Document emergency bank and payroll continuity, corporate access and passwords held securely, responsible board contacts, customer/learner obligations, critical contracts, legal notices, short-term CEO delegation, key-person insurance review where sensible, and a plan for secure access by designated authorized persons. Review annually and after major funding.

\section{Official Sources and Research Foundations}

The following statutory, regulatory, and institutional references informed the principles of this chapter:

\begin{enumerate}
    \item \textbf{[G1] Ministry of Corporate Affairs: SPICe+ and Incorporation FAQs.} Statutory guidelines and incorporation context for Section 8 and commercial entities. \url{https://www.mca.gov.in/content/dam/mca/pdf/SPICEplus-and-linked-filings-FAQs-V3-20230122.pdf}
    \item \textbf{[G2] Ministry of Corporate Affairs: Companies Act, 2013.} Primary legislation, director duties, related-party transactions, securities issuance, and accounts. \url{https://www.mca.gov.in/Ministry/pdf/CompaniesAct2013.pdf}
    \item \textbf{[G3] Ministry of Electronics and Information Technology: DPDP Rules 2025 and Enforcement Timeline.} Digital Personal Data Protection regulatory rules, consent architecture, and data fiduciary compliance. \url{https://www.meity.gov.in/documents/act-and-policies/digital-personal-data-protection-rules-2025-gDOxUjMtQWa?pageTitle=Digital-Personal-Data-Protection-Rules-2025}
    \item \textbf{[G4] Reserve Bank of India: Foreign Investment Reporting Regulations and FIRMS Portal.} FEMA notifications, FC-GPR, FC-TRS, and annual foreign liabilities and assets reporting. \url{https://rbi.org.in/SCRIPTS/BS_FemaNotifications.aspx?Id=11723} and \url{https://firms.rbi.org.in/firms/faces/pages/login.xhtml}
    \item \textbf{[G5] Ministry of Home Affairs: Foreign Contribution Regulation Act (FCRA) Portal.} Statutory rules distinguishing ordinary commercial fee income from foreign donations and grants. \url{https://fcraonline.nic.in/}
    \item \textbf{[G6] University Grants Commission (UGC): Statutory Guidelines and FAQs.} Official definitions of lawful degree-awarding authority and state private university standards. \url{https://www.ugc.gov.in/Home/faq}
    \item \textbf{[G7] Maharashtra Higher and Technical Education Department: Private Universities Act and 2024 Rules.} Legal framework governing state private university sponsorship, detailed project reports (DPR), and statutory academic councils. \url{https://htedu.maharashtra.gov.in/Main/MaharashtraPrivateUniversitiesAct}
    \item \textbf{[G8] World Intellectual Property Organization (WIPO): Institutional IP Policy Toolkit.} Models and frameworks for university and corporate technology transfer, invention disclosure, and licensing. \url{https://www.wipo.int/en/web/technology-transfer/ip-policies}
    \item \textbf{[G9] OECD (2015, updated frameworks): \textit{G20/OECD Principles of Corporate Governance}.} Benchmark standards for shareholder rights, corporate disclosure, board oversight, and stakeholder protections.
    \item \textbf{[G10] Robert A. G. Monks and Nell Minow: \textit{Corporate Governance}.} Foundational treatise on corporate agency, fiduciary duties, board accountability, and shareholder rights.
    \item \textbf{[G11] David F. Larcker and Brian Tayan: \textit{Corporate Governance Matters}.} Empirical analysis of board dynamics, executive compensation, monitoring, and startup governance.
    \item \textbf{[G12] The Bayes Institute Reference: \textit{The Theory of Institutional Evolution} (2026).} Analytical framework separating commercial operating enterprises from research institutes and academic bodies across the 3 capital markets.
\end{enumerate}

\subsection*{Legal Source Caution}

Regulatory websites can contain historical versions; binding obligations derive from operative law, notifications, rules and the specific transaction. This chapter explains what to ask and document, not what a particular company has already complied with. Obtain independent legal opinions for financing, entity restructuring, property, foreign funds, and degree authority.

\section*{Closing Principle}

\begin{quote}
\itshape
\textbf{Governance is the operating constitution that preserves agency.} Governance is not an abstract virtue reserved for mature corporations or ancient universities. It is the real-time machinery that determines who can commit the company, whose capital is at risk, what information can be trusted, who can halt a fatal error, and how founders preserve legitimate strategic control. The enduring institution separates commercial returns from academic truth, documents every right before conflict emerges, and \textbf{builds a decision system capable of outlasting its founders.}
\end{quote}
